
\renewcommand{\bibname}{\textbf{\Huge\textcolor{bleu}{Bibliographie}}}
\addcontentsline{toc}{chapter}{Bibliographie}
\begin{thebibliography}{40}

\bibitem{ref1} I. N. Sonatrach, « Historique et évolution », Institut numérique, 2012. [En ligne]. Disponible : \url{https://www.institut-numerique.org/2-historique-et-evolution-5028f58cb09d3}

\bibitem{ref2} J. Schneider, S. Susnjara et I. Smalley, « Hyperconverged infrastructure explained », IBM Think, févr. 2024. [En ligne]. Disponible : \url{https://www.ibm.com/think/topics/hyperconverged-infrastructure}

\bibitem{ref4} A. Pande, « Proxmox vs. Harvester: Can the enterprise-grade platform beat the community favorite? », XDA Developers, févr. 2025. [En ligne]. Disponible : \url{https://www.xda-developers.com/proxmox-vs-harvester/}

\bibitem{ref5} SUSE, « Harvester HCI Documentation », Harvester HCI. [En ligne]. Disponible : \url{https://harvesterhci.io/}

\bibitem{ref6} Proxmox Server Solutions, « Proxmox Virtual Environment Overview », Proxmox. [En ligne]. Disponible : \url{https://www.proxmox.com/en/products/proxmox-virtual-environment/overview}

\bibitem{ref8} Open vSwitch Authors, « What is Open vSwitch? », Documentation Open vSwitch. [En ligne]. Disponible : \url{https://docs.openvswitch.org/en/latest/intro/what-is-ovs/}

\bibitem{ref13} Red Hat, « Chapter 1: The Ceph Architecture », Red Hat Ceph Storage 4 Architecture Guide. [En ligne]. Disponible : \url{https://docs.redhat.com/en/documentation/red_hat_ceph_storage/4/html/architecture_guide/the-ceph-architecture_arch}

\bibitem{ref14} A. Chandak, K. Jaju, A. Kanfade, P. Lohiya et A. Joshi, « Dynamic Load Balancing of Virtual Machines using QEMU-KVM », \textit{International Journal of Computer Applications}, vol.~46, n\textsuperscript{o}~6, pp.~10--14, mai 2012.

\bibitem{ref17} J. Watts, « A practical approach to dynamic load balancing », Mémoire de master, California Institute of Technology, 1995. [En ligne]. Disponible : \url{https://thesis.caltech.edu/6924/1/Watts_j_1995.pdf}

\bibitem{ref20} Corosync Project, « The Corosync Cluster Engine », Corosync. [En ligne]. Disponible : \url{https://corosync.github.io/corosync/}

\bibitem{ref25} Grafana Labs, « Grafana Documentation », Grafana Labs. [En ligne]. Disponible : \url{https://grafana.com/docs/grafana/latest/}

\bibitem{ref26} W. Tarreau, « HAProxy version 2.8: configuration manuelle », HAProxy. [En ligne]. Disponible : \url{https://www.haproxy.org/download/2.8/doc/configuration.txt}

\bibitem{ref27} S.~A. Weil, S.~A. Brandt, E.~L. Miller et D.~D. Long, « Ceph: A scalable, high-performance distributed file system », dans \textit{Proceedings of the 7th USENIX Symposium on Operating Systems Design and Implementation (OSDI'06)}, USENIX Association, 2006, pp.~307--320.

\bibitem{ref28} S.~A. Weil, S.~A. Brandt, E.~L. Miller et C. Maltzahn, « CRUSH: Controlled, Scalable, Decentralized Placement of Replicated Data », dans \textit{Proceedings of the 2007 ACM/IEEE Conference on Supercomputing (SC'07)}, ACM, 2007.

\bibitem{ref33} Proxmox Server Solutions, « Proxmox VE Ceph Server », Proxmox VE Wiki, 2024. [En ligne]. Disponible : \url{https://pve.proxmox.com/wiki/Ceph_Server}

\bibitem{ref38} T.~C. Wilcox~Jr., « Dynamic Load Balancing of Virtual Machines Hosted on Xen », Mémoire de master, Brigham Young University, avr. 2009. [En ligne]. Disponible : \url{https://scholarsarchive.byu.edu/cgi/viewcontent.cgi?article=2653&context=etd}

\bibitem{ref40} GitLab, « GitLab Requirements », Documentation GitLab. [En ligne]. Disponible : \url{https://docs.gitlab.com/install/requirements/}

\end{thebibliography}
